\section{Introduction}

An agentic coding assistant pointed at an HPC codebase does what such agents do: it writes a parallel kernel, compiles it, runs it, compares the output against a trusted reference, and - seeing them agree - reports the task finished. That self-check is how modern agents decide their own work is done, and increasingly it is also how the reward that post-trains them is computed. For most software it is a sound rule. For high-performance code it is blind in the one direction that matters, because the thing the code exists to deliver - speed across many cores - never appears in the output it checks. An autonomous HPC agent judged this way will accept, and a training loop rewarded this way will reinforce, a program that is perfectly correct and slower than the sequential version it was meant to replace. The agent does not flag it, because nothing it measures can see it. This is the problem we take on: not that agents write buggy parallel code, but that ``correct'' has become easy and ``fast'' is invisible to the loop that decides when an agent is done and trains the next one.

The checker became the objective. Code models are increasingly post-trained not on human preference but on execution: the model writes a program, an automatic checker runs it, and the reward is paid only if the program does what it was asked to do~\cite{tulu3,deepseekr1,grpo}. The same substitution has happened inside agents, which decide for themselves that a draft is finished by running it and comparing its output to what was expected. In both settings the checker has moved out of test infrastructure and into the objective function. Whatever it scores highly is what gets written - so if the checker looks at the wrong thing, so, eventually, does the model.

For parallel code the checker is the one everyone already has: compile the generated program, run it, compare its output against a trusted serial reference, and pay the reward if they agree. This is what the standard benchmark for LLM-generated parallel code does~\cite{pareval}, what an execute-and-fix agent does when it accepts its own OpenMP draft~\cite{selfdebug,reflexion}, and the natural verifiable reward to reach for when post-training an HPC code model. It is also what the current wave of agentic HPC systems quietly rely on beneath their search - self-correction loops, multi-agent orchestration, auto-tuning, evolutionary and tree search - however elaborate the machinery stacked on top. It is the obvious checker. For parallel code it is also the wrong one, and not for want of tuning: because the flaw is in the signal every one of those systems shares, more search cannot remove it, and the fix has to live in the signal itself.

The reason is a mismatch between what the checker reads and what the code is for. Correctness is a predicate on \emph{outputs}: it asks whether the program computed the right numbers. But the entire reason to write parallel code rather than a sequential loop is a property of \emph{time} - that it finishes sooner when handed more cores. Two programs that return identical results, one that scales across eight cores and one that is slower than a plain sequential loop, are indistinguishable to a checker that reads only outputs. The reward is therefore blind, by construction, to the one property the code was written to have. More testing does not close this gap: additional tests sharpen the output check, and the output check is not where the missing information lives.

Two consequences follow, and they differ in kind. The first concerns the \emph{optimum}: nothing in a correctness-only reward pays for parallelism, so no procedure that maximizes it has any reason to emit parallel code. Where models do write it under such an objective they are spending a habit inherited from pretraining, which the reward neither reinforces nor protects and can erode at no cost. The second concerns the \emph{gradient}, and it bites long before the optimum is approached. Group-relative estimators such as GRPO center each sampled group's reward on its mean and divide by its standard deviation; when correctness is near-saturated - \PassRate\% of our single-shot generations are already correct at every thread count we test - a sampled group comes back uniformly correct, its reward is uniformly $1$, its variance is zero, and the advantage vanishes. The reward supplies no signal, not because the model is right, but because it has nothing further to say. A sharper correctness checker makes this worse: correctness is the one axis on which there is nothing left to win.

We turn both claims into measurements on a frozen pool of \NAccepted\ verified programs, so that every number in this paper can be recomputed offline without an inference budget, and we build the system the argument demands: a grounded diagnostician that answers questions about a generated program only by calling real measurement tools - compile, sweep thread counts, time against the serial reference, check for races - paired with a \emph{performance-critic} that refuses any claim a measurement does not support (Section~\ref{sec:system}). The efficiency gate that critic enforces at inference is exactly the reward we recommend for training: the checker and the reward are one object, and fixing one fixes the other.

The whole argument fits in a single experiment, and we found it by trying to break our own verifier. Take a correct OpenMP program and delete every \texttt{\#pragma omp} line in it. It still compiles. It still runs. It still agrees with the serial reference on every input at every thread count - \emph{more} reliably than before, in fact, because there is now nothing left to race. So it still earns the maximum reward, and the serial program sits in the argmax. This is not a corner case to be patched: the reward is a predicate on outputs, deleting the directives preserves the outputs by construction, and so no amount of additional testing can remove it. When we ran exactly this transform across the pool, \SPRate\% of the programs our verifier had accepted stayed correct at a mean reward change of exactly \SPRewardDeltaCorr - \NAccepted\ instances, not an argument. A reward for parallel code that is maximized by deleting the parallelism is not measuring parallel code, and the rest of this paper is what follows once you take that seriously.

\paragraph*{What we do not claim}
The literature has moved quickly and we want the boundary of the contribution to be exact. We are \emph{not} the first to observe that model-generated parallel code can be functionally correct while executing serially: PARACODEX~\cite{paracodex} documents this as ``bypass'' and ``pseudo-offload'' in a GPU setting and characterizes it as an agent and harness failure to be patched with harness constraints. Our difference is what the observation is a fact \emph{about}-not a bug in a particular harness, but the location of the optimum of the reward function itself, on CPU thread scaling rather than device offload. We are \emph{not} the first to combine correctness with measured performance in a reward for code: ACECode~\cite{acecode} does so additively for Python runtime; Kevin~\cite{kevin}, CUDA-L1~\cite{cudal1}, CUDAPerf~\cite{cudaperf} and MaxCode~\cite{maxcode} do so for GPU kernels; RLPF~\cite{rlpf} aligns code models to performance through a learned reward model and reports 1.9--4.5$\times$ on OpenMP; and real-machine GFLOPS rewards have been used to post-train an HPC code model with GRPO~\cite{realmachine}. Our contribution is not the combination, but the argument for why the correctness term cannot stand alone, the structure that follows from it-a multiplicative gate rather than a weighted sum-and an adversarial suite that lets these formulations be compared rather than merely listed. Finally, we are \emph{not} claiming that best-of-$N$ selection is equivalent to group-relative policy optimization: it is a proxy for the selection component only~\cite{bonmaxk}, since GRPO additionally sharpens the sampling distribution~\cite{rewardunlikely}, and we label our selection analysis as such throughout.

\paragraph*{Contributions}
\begin{enumerate}
  \item \textbf{Optimum degeneracy.} A statement and proof, via the serial projection, that correctness-only verifiable rewards for parallel code are maximized by programs with no parallelism, with the corollary that group-relative estimators suffer \emph{signal collapse} on saturated correctness distributions (Section~\ref{sec:degeneracy}).
  \item \textbf{An executable demonstration.} Pragma-stripping the accepted pool preserves correctness for \SPRate\% of programs at zero reward cost, turning the argument into instances (Section~\ref{sec:results-projection}).
  \item \textbf{A quantification of the blind spot.} Accepted programs range from below-serial to near-linear scaling at identical reward---a spread correctness-only selection cannot navigate, because it ties every candidate and must pick blindly (Sections~\ref{sec:results-contamination}, \ref{sec:results-bakeoff}).
  \item \textbf{\PG.} An adversarial reward-hack suite for parallel code, and a bake-off of six reward formulations against it-including the cases our own gate fails (Section~\ref{sec:results-hacks}).
  \item \textbf{A measurement protocol and a released artifact.} A timing protocol that closes the measurement-corruption attacks (fresh processes, whole-driver timing, randomized inputs), together with an artifact that recomputes every reward and figure in the paper offline (Section~\ref{sec:results-robust}).
\end{enumerate}